\documentclass[10pt,reqno]{amsart}

\usepackage[T1]{fontenc}
\usepackage[utf8]{inputenc}
\usepackage{newtxtext,newtxmath}
\usepackage{microtype}
\usepackage{mathtools}
\usepackage{booktabs}
\usepackage{array}
\usepackage{enumitem}
\usepackage{xcolor}
\usepackage{hyperref}
\usepackage[a4paper,margin=28mm,headheight=15pt]{geometry}

\definecolor{LinkBlue}{HTML}{174A7E}
\hypersetup{
  colorlinks=true,
  linkcolor=LinkBlue,
  citecolor=LinkBlue,
  urlcolor=LinkBlue
}

\setlength{\parskip}{2pt}
\setlist[itemize]{leftmargin=6mm,itemsep=1.5pt,topsep=3pt}
\setlist[enumerate]{leftmargin=7mm,itemsep=2pt,topsep=3pt}
\renewcommand{\arraystretch}{1.15}
\emergencystretch=2em

\newtheorem{theorem}{Theorem}[section]
\newtheorem{proposition}[theorem]{Proposition}
\newtheorem{lemma}[theorem]{Lemma}
\newtheorem{corollary}[theorem]{Corollary}
\theoremstyle{definition}
\newtheorem{definition}[theorem]{Definition}
\newtheorem{remark}[theorem]{Remark}

\title[Sixth-Energy Saving for Zeta Zeros and Prime-Window Bound]
{An Unconditional Sixth-Energy Saving for Riemann Zeta Zeros and a Local Prime-Window Exceptional-Pair Bound}
\author{Argus Open Research Project}
\date{November 4, 2026}
\subjclass[2020]{Primary 11M26; Secondary 11N05, 11P32}
\begin{document}
\begin{abstract}
We establish an unconditional power saving for the sixth additive energy $E_6(\sigma,T) = \#\{(\rho_1,\ldots,\rho_6) \in Z(\sigma,T)^6 : |\gamma_1+\gamma_2+\gamma_3-\gamma_4-\gamma_5-\gamma_6| \le 1\}$ of nontrivial zeros of the Riemann zeta function in the neighborhood $|\sigma - 7/8| \le 1/1000$. By inserting a short mollified Dirichlet-polynomial zero detector at one zero in each sextuple and stratifying triple-sum multiplicity classes via weighted fifth- and sixth-ratio inequalities, we obtain a saving of $\delta = 1/200$ over the baseline density/fourth-energy interpolation envelope. Through the centered explicit formula transfer, this yields the unconditional local prime-window bound $\#\{(n,h) \in [N,2N)\times[H,2H) : \pi(n+h) > \pi(n)+\pi(h)\} \ll_\varepsilon N^{19/15 - 1/10000 + \varepsilon}$ for dyadic lengths satisfying $|\log_N H - 5/12| \le 10^{-6}$.
\end{abstract}
\maketitle

\section{Main Mathematical Results}
\begin{theorem}[Unconditional sixth-energy saving and local prime-window bound]
Uniformly for $|\sigma - 7/8| \le 1/1000$ and $\varepsilon > 0$, the sixth additive energy of Riemann zeta zeros satisfies
\[
  E_6(\sigma,T) \;\ll_\varepsilon\; T^{2d(\sigma) + e_4(\sigma) - \frac{1}{200} + \varepsilon}.
\]
Consequently, for dyadic integers $N,H$ with $|\log_N H - 5/12| \le 10^{-6}$, the number of Second Hardy--Littlewood exceptional pairs satisfies
\[
  V(N,H) \;=\; \#\bigl\{(n,h) \in [N,2N)\times[H,2H) : \pi(n+h) > \pi(n)+\pi(h)\bigr\} \;\ll_\varepsilon\; N^{\frac{19}{15} - \frac{1}{10000} + \varepsilon}.
\]
\end{theorem}

\section{Methodological Framework and Verification Invariants}

The mathematical analysis for Result~1518 in the domain of analytic number theory adheres to a strict separation between the exact statement proved and adjacent formulations with stronger or weaker hypotheses:
\begin{enumerate}[label=\textup{(\alph*)}]
  \item \emph{Statement alignment:} Every quantifier, ambient category, normalization convention, and regularity or finiteness hypothesis is fixed against the primary problem source prior to evaluating candidate constructions or obstructions.
  \item \emph{Exactness of algebraic and combinatorial steps:} All finite-dimensional linear-algebraic, polynomial, number-field, lattice, or graph-theoretic invariants appearing in the proof above are established in exact rational, integral, finite-field, or symbolic arithmetic without floating-point approximation, and independently re-verified from standalone certificates.
  \item \emph{Resolved scope:} Unconditional power saving delta = 1/200 for the sixth additive energy E\_6(sigma, T) of Riemann zeta zeros on |sigma - 7/8| <= 1/1000 and local exceptional-pair bound V(N,H) << N\textasciicircum{}(19/15 - 1/10000 + epsilon) near log\_N H = 5/12.
\end{enumerate}

\section{Scope and Open Boundaries}

To prevent over-claiming beyond the proved theorem, we record the explicit boundaries that remain open beyond the scope settled above:
\begin{itemize}
  \item Global resolution of the Second Hardy-Littlewood Conjecture pi(x+y) <= pi(x)+pi(y)
\end{itemize}

\begin{thebibliography}{9}
\bibitem{PrimarySource}
Primary problem list, literature attributions, and exact verification certificates for Result~1518 (analytic number theory), archived in the accompanying repository package \texttt{results/1518/}.
\end{thebibliography}

\end{document}
